\section{Natural root diagrams and Cartier transfers}
\label{sec:transfers}

The invariant-function calculation makes the source comparison unit
below constant and lets us normalize a single Cartier transfer on a
finite marked algebra. Compact-support finiteness enters the Laurent
argument in \Cref{sec:laurent}. The construction must work on entire finite
diagrams, and its pole bound must survive every iterate: these are the
two inputs the later pro-transfer argument will need.

Fix an exact-height stratum with $1\leq m<n$ and put $r=n-m$.
Throughout this section $G=P_n$ is unchanged.  A sufficiently deep product
subgroup $U\subset P_m\times\GL_r(\Zp)$ will be chosen in the course of the
construction.  Write
\[
 A=\kk[[x,y_1,\ldots,y_{r-1}]],\qquad B=A[1/x],\qquad
 R=B_U=eB',\qquad B'=A'[1/x]
\]
as in \Cref{prop:mark-whole-model}.  In particular, $A'$ is the whole finite
free $A$-algebra with its fixed integral basis, $B'/B$ is finite \'{e}tale,
and $e$ is a $G$-invariant idempotent.  No specialization of $e$ across
$x=0$ will be used.

\subsection{The distribution algebra and the root functor}

Let $B_\infty=\colim_U B_U$ denote the algebra with both full markings.
The compatible lifts of the marked \'{e}tale generators form a torsor
$\mathscr P$ under the affine group scheme
\[
 T=\varprojlim_{[p]}\Gamma_m[p^j]^r.
\]
These are affine torsors, understood through their finite faithfully flat
quotients.  By \Cref{prop:mark-roots},
\begin{equation}
 \cO(\mathscr P)=\colim_j B_\infty^{1/p^{mj}}
                 =B_\infty^{\mathrm{perf}},
 \qquad
 \mathscr P/[p^\ell]T=\Spec B_\infty^{1/p^{m\ell}}.
 \label{eq:trans-full-torsor}
\end{equation}
The $G$-action commutes with $T$; changes of markings act on this identity
and on $T$ by the indicated constant Honda automorphisms.

Here and below a root of a module has a specific meaning.  If $h$ is a
power of $p$ fixing $\kk$ and $M$ is an $R$-module, $M^{1/h}$ is the copy of
$M$ transported along the isomorphism
\[
 R\longrightarrow R^{1/h},\qquad b\longmapsto b^{1/h}.
\]
It is then regarded as an $R$-module by the \emph{inclusion}
$R\subset R^{1/h}$.  Thus a finite locally free $M$ of rank $v$ gives an
$R$-module $M^{1/h}$ of rank $h^rv$.  The map $M^{1/h}\to M$ denoted by
$w\mapsto w^h$ is the transport isomorphism; it is semilinear with respect
to $R^{1/h}\to R$, $b\mapsto b^h$.  This convention applies to maps and
complexes, and does not mean pullback along $R\subset R^{1/h}$ followed
by retaining its rank over the larger ring.

For the remainder of this subsection assume $m>1$. The height-one
case is treated separately in the transfer construction below.
The continuous distribution algebra is
\begin{equation}
 \Lambda=\cO(T)^\vee\simeq\kk[[D_1,\ldots,D_d]],
 \qquad d=(m-1)r.
 \label{eq:trans-distributions}
\end{equation}
Here the dual is the inverse limit of the duals of the finite Hopf
algebras.  We recall why both the dimension and the powers in this
description have these values.  Cartier duality identifies this inverse
limit with the coordinate algebra of the formal group
$(\Gamma_m^D)^r$.  The standard facts used here are the equivalence
between connected $p$-divisible groups and formal Lie groups for which
$[p]$ is finite faithfully flat
over a complete local Noetherian base of residue characteristic $p$,
and finite-flat Cartier duality, which interchanges Frobenius and
Verschiebung; see \citet[\S1.2 and \S\S2.2--2.3, Propositions~1 and~3]{tate1967}.
Our base is the perfect field $\kk$.  On the one-dimensional Honda
group, $F^m=[p]$ and $FV=[p]$.  Duality gives $V^m=[p]$ on the dual;
there $F^m[p]=[p]^m$, and cancellation of the faithfully flat isogeny
$[p]$ gives
\begin{equation}
 F^m=[p]^{m-1}\quad\hbox{on }\Gamma_m^D.
 \label{eq:trans-dual-honda}
\end{equation}
All Frobenius twists in this formula use the Honda model over $\Fp$.
The kernel of $F$ on $\Gamma_m$ has order $p$, whereas $\Gamma_m[p]$ has
order $p^m$.  Thus the kernel of $V$ has order $p^{m-1}$, and the dual
Frobenius has that degree.  A smooth formal group of dimension $s$ has
Frobenius degree $p^s$, so $\dim\Gamma_m^D=m-1$.  When $m>1$ the dual has
no \'{e}tale part: on any finite \'{e}tale $p$-power subgroup Frobenius is
invertible, whereas a sufficiently high iterate of the right side of
\cref{eq:trans-dual-honda} is zero.  The connected $p$-divisible group is
therefore formal, giving \cref{eq:trans-distributions}.
These uses of Cartier duality are at finite flat levels before passing
to the limit.

Choose a positive integer $a_0$ divisible enough to fix the finite
constant field, and set, for $a\geq1$,
\begin{equation}
 q_0=p^{ma_0},\quad q=q_0^a,\quad
 \ell=(m-1)a_0a,\quad h=q^{m-1}=p^{m\ell}.
 \label{eq:trans-powers}
\end{equation}
On the dual formal group, \cref{eq:trans-dual-honda} identifies
$[p^\ell]$ with the $q$-power Frobenius.  Its map on coordinates is the
inclusion
\[
 \sigma_q:\Lambda\longrightarrow\Lambda,
 \qquad D_i\longmapsto D_i^q.
\]
Write $\Lambda^{(q)}=\sigma_q(\Lambda)$ and
\[
 \Lambda_q=\Lambda/(D_1^q,\ldots,D_d^q),\qquad
 \mathsf N_q=\Hom_\kk(\Lambda_q,\kk).
\]
The quotient $T/[p^\ell]T$ has coordinate representation $\mathsf N_q$.
In particular
\[
 \dim_\kk\mathsf N_q=q^d=h^r=p^{m\ell r},
\]
in agreement with the rank of the finite lift torsor.

For a finite-length $\Lambda$-module $E$, regard $E$ as a
$\Lambda^{(q)}$-module through $\sigma_q^{-1}$ and define
\[
 \Coind_q(E)=\Hom_{\Lambda^{(q)}}(\Lambda,E),
 \qquad (a\cdot f)(b)=f(ab).
\]
This is an exact functor, because $\Lambda$ is free over
$\Lambda^{(q)}$ on the monomials $D^\alpha$, $0\leq\alpha_i<q$.
For $E=\kk$ it gives $\mathsf N_q$.  We use the natural change-of-markings
action on coinduction.  Explicitly, if $u$ acts compatibly on $E$, its
action is $(u f)(b)=u(f(u^{-1}b))$.  The inclusion $\sigma_q$ commutes
with this action: its coefficients lie in the finite field fixed by $q$.

Denote by $V_E$ the associated coefficient bundle, obtained by descent
of $\mathscr P\times E$ through translations and then through $U$.
Thus $V_\kk=R$.  A fixed finite diagram of such representations and maps
descends to finite locally free bundles at some finite marking stage.
A plain finite $\Lambda$-module diagram can be given trivial auxiliary
action after shrinking $U$: it factors through a finite-dimensional
quotient of $\Lambda$, whose marking action has finite image.
Coinduced diagrams of arbitrarily large size retain their natural
auxiliary actions; no single $U$ is asserted to act trivially on all of
them.

\begin{proposition}[Natural roots of finite diagrams]
\label[proposition]{prop:trans-roots}
Assume $m>1$. With the powers in \cref{eq:trans-powers}, there is a natural
$G$-equivariant isomorphism
\begin{equation}
 V_{\Coind_q(E)}\simeq (V_E)^{1/h}.
 \label{eq:trans-root-functor}
\end{equation}
It is compatible with the natural auxiliary actions and with every
arrow of a finite diagram.  It preserves exact diagrams and is
compatible with successive coinduction and successive roots.  For any
fixed initial diagram, these assertions hold at every sufficiently
deep finite marking stage, retaining the natural auxiliary action on
each of its coinductions.
\end{proposition}

\begin{proof}
The perfectness in \cref{eq:trans-full-torsor} makes the $h$-power map
an isomorphism on the \emph{total lift torsor}.  Write
$\rho:\cO(\mathscr P)\to\cO(\mathscr P)\otimes_\kk\cO(T)$ for its
coaction.  Absolute powering satisfies
\begin{equation}
 \rho(f^h)=\rho(f)^h
     =(\Fr_h^*\otimes[p^\ell]^*)\rho(f),
 \qquad f\in\cO(\mathscr P).
 \label{eq:trans-power-coaction}
\end{equation}
Here $\Fr_h^*$ raises \emph{all} coefficients, including the base
parameters, to the $h$th power.  Thus this is a map over the matching
Frobenius base, not an endomorphism over the original marked base.
The Honda relation identifies the translation Frobenius with
$[p^\ell]$.  In torsor-action notation the same identity is
$\Fr_h(t\cdot z)=[p^\ell]t\cdot\Fr_h(z)$.
This endomorphism of $T$ gives an isomorphism
\[
 T\xrightarrow{\sim}[p^\ell]T
\]
onto a subgroup of $T$.  It is not an automorphism onto all of $T$.
For example the projection of $[p^\ell]T$ to the level-$\ell$ quotient
is zero.  The group $T$ need not be reduced, although the total torsor
in \cref{eq:trans-full-torsor} is perfect.

It follows from \cref{eq:trans-power-coaction} that, over
$B_\infty^{1/h}$, the same total lift space with structure group
$[p^\ell]T$ is the root of the original $T$-torsor.  Associated-bundle
induction in stages now gives
\[
 V_{\Coind_{[p^\ell]T}^{T}(E)}
   =\pi_*\bigl(V_E\text{ on }\mathscr P/[p^\ell]T\bigr)
   \simeq (V_E)^{1/h},
\]
where $\pi:\Spec B_\infty^{1/h}\to\Spec B_\infty$ is the finite root
map.  The pushforward is the restriction of scalars specified above.
One can verify the equality without a choice of a torsor point:
after a faithfully flat trivialization it is the usual evaluation
isomorphism for coinduced equivariant functions, and that isomorphism
commutes with the two projections in the descent diagram.  It therefore
descends, is natural in $E$, and identifies each morphism of any finite
diagram.  This also proves exactness, since the trivialized functors
are exact and the base changes are faithfully flat.

All the maps used here commute with $G$.  They also commute with the
changes of markings, since $[p^\ell]$ is central and the chosen power
fixes $\kk$.  For two powers, the identities
$(z^{h_1})^{h_2}=z^{h_1h_2}$ and
$[p^{\ell_2}][p^{\ell_1}]=[p^{\ell_1+\ell_2}]$ show that the two
resulting torsor maps are equal.  The evaluation map for coinduction
in stages has the same composition law.  This proves the asserted
coherence, including the maps of the diagrams, rather than only an
identification of their terms.

Finally, the original finite diagram, its torsor quotient, and its
coaction involve finitely many matrix coefficients and relations.
They descend through the filtered union of finite \'{e}tale marking
algebras.  Choose one stage containing this data and shrink $U$ so that
the original diagram has the required auxiliary action there.  At
every coinduction level the preceding identity is still the same
intrinsic power identity, with its natural auxiliary action; it
therefore descends to that stage.  Equivariance can be checked after
the one faithfully flat full-marking extension, since $G$ preserves
each finite marking algebra.  This does not require intersecting
depth conditions for individual elements of $G$, or making all high
coinduced representations trivial under $U$.
\end{proof}

The elementary pairing underlying coinduction will also be useful.
Let $\kappa_q:\Lambda\to\Lambda^{(q)}$ extract the coefficient of
$\prod_iD_i^{q-1}$ in the displayed free basis.  The pairing
$(a,b)\mapsto\kappa_q(ab)$ is perfect: modulo the maximal ideal of
$\Lambda^{(q)}$ its matrix pairs $D^\alpha$ with
$D^{(q-1)\mathbf1-\alpha}$, and hence is invertible; its determinant
is a unit already over $\Lambda^{(q)}$.  Thus coinduction is naturally
isomorphic to scalar extension as a functor of plain modules, once
this functional is fixed.  This last choice need not identify the
auxiliary actions, so the action supplied by
\Cref{prop:trans-roots} will always be retained.

\subsection{A volume at one finite marking stage}

Use continuous differentials on the formal parameter ring and their
localizations and finite \'{e}tale base changes; write this module as
$\widehat\Omega^1_{R/\kk}$.  It is free of rank $r$ before any perfection.

For a finite translation quotient $H$, a \emph{translation augmentation}
means a nonzero $H$-equivariant map $\lambda:\cO(H)\to\kk$, with trivial
action on $\kk$. Equivalently, it is an invariant integral:
$(\id\otimes\lambda)\Delta(a)=\lambda(a)1$.
For $m>1$ and the quotient indexed by $q$, this is a $\Lambda$-linear
map $\mathsf N_q\to\kk$. The term refers to this invariant functional,
not to the Hopf counit.

\begin{proposition}[Invariant volume and scalar comparison]
\label[proposition]{prop:trans-volume}
After one finite marking shrink to a pro-$p$ subgroup,
$\widehat\Omega^1_{R/\kk}$ has a
$G$-invariant frame.  In particular it has a $G$-invariant volume
generator $\eta\in\widehat\Omega^r_{R/\kk}$.

For a finite root order $h$ fixing $\kk$, the normalized Cartier
functional
\begin{equation}
 T_{\eta,h}:R^{1/h}\longrightarrow R,
 \qquad T_{\eta,h}(z)=\frac{\Cartier^{\log_p h}(z^h\eta)}{\eta}
 \label{eq:trans-cartier-functional}
\end{equation}
is $R$-linear and generates $\Hom_R(R^{1/h},R)$ as a module over
$R^{1/h}$.  A nonzero translation augmentation at the corresponding
finite torsor level induces a scalar multiple of this functional;
the scalar belongs to the chosen finite constant field.
\end{proposition}

\begin{proof}
The intermediate order-$p$ root cover is a torsor under the height-one
Frobenius quotient of translations, of rank $p^r$.  Over the chosen
Honda constants this quotient is $\alpha_p^r$ for $m>1$ and $\mu_p^r$
for $m=1$.  This follows by taking the first Frobenius kernel of the
Honda group; its Verschiebung is zero in the first case and invertible
in the second.  These finite Hopf algebras and their Lie spaces are
defined over a finite field.  Their continuous marking actions factor
through a finite group.  Shrink $U$ to fix the quotient and a Lie basis.
For $m=1$ choose the standard toral basis of the fixed $\mu_p^r$.
Finite presentation of the coaction descends its torsor action to a
finite stage.  Torsor freeness identifies the resulting invariant
vector fields with a basis
\[
 \delta_1,\ldots,\delta_r
 \quad\hbox{of }\operatorname{Der}_R(R^{1/p},R^{1/p}).
\]
For completeness this freeness can be checked after trivializing the
torsor: the vectors are $\partial/\partial z_i$ on $\alpha_p^r$ and
$z_i\partial/\partial z_i$ on $\mu_p^r$.  Their coefficient matrices
are invertible, and invertibility descends faithfully flatly.

For $b\in R$, define
\begin{equation}
 \partial_i(b)=\bigl(\delta_i(b^{1/p})\bigr)^p.
 \label{eq:trans-frame}
\end{equation}
This formula is additive, satisfies the ordinary Leibniz rule, and is
$\kk$-linear because $\kk$ is perfect.  To see that it is a frame, let
$(t_1,\ldots,t_r)=(x,y_1,\ldots,y_{r-1})$, a $p$-basis after the
finite \'{e}tale base change.  Write
\[
 \delta_i=\sum_j a_{ij}\frac{\partial}{\partial t_j^{1/p}}.
\]
Then $\partial_i=\sum_j a_{ij}^p\partial/\partial t_j$ and the new
determinant is $\det(a_{ij})^p\in R^\times$.  The derivations commute
with $G$, because the full-marking action commutes with translations
and with powers.  The identities descend to this same finite stage:
they hold after the faithful full-marking extension, and $G$ preserves
$R$.  The dual frame and its top wedge give $\eta$.  All differentials
in this construction were taken at a finite root stage.

We describe the Cartier operator used in
\cref{eq:trans-cartier-functional}.  If
$dt=dt_1\wedge\cdots\wedge dt_r$ and $h=p^\nu$, the unique $p$-basis
expansion of a form is
\[
 f\,dt=\sum_{0\leq\alpha_i<h} f_\alpha^h t^\alpha\,dt,
 \qquad f_\alpha\in R.
\]
The operator is
\begin{equation}
 \Cartier^\nu(f\,dt)=f_{(h-1)\mathbf1}\,dt.
 \label{eq:trans-cartier-coordinates}
\end{equation}
For power series the expansion groups the exponents modulo $h$;
localization and finite \'{e}tale base change preserve it.  First take
$h=p$.  In one variable the de Rham differential on a residue exponent
other than $p-1$ gives an exact term, while the last exponent gives the
class of $t^{p-1}dt$.  Tensoring these one-variable complexes proves
the Cartier isomorphism in $r$ variables, and iterating it gives
\cref{eq:trans-cartier-coordinates} for $h=p^\nu$.  The inverse of the
one-step isomorphism is
characterized on one-forms by $db\mapsto[b^{p-1}db]$ and by
multiplication on exterior powers.  This characterization shows that
the formula commutes with coordinate changes and \'{e}tale base change.
For the customary inverse-Frobenius convention and its duality
interpretation, see \citet[Definition~2.1 and Example~2.23]{blickleBockle2011}.
It is continuous on the indicated formal coefficient modules.  In
particular it commutes with $G$ and satisfies
$\Cartier^\nu(b^h\omega)=b\Cartier^\nu(\omega)$.

\Cref{eq:trans-cartier-coordinates} proves the asserted
linearity and dual-generator property.  More explicitly, for
$\eta=dt$ the pairing
\[
 (u,v)\longmapsto T_{dt,h}(uv)
\]
on $R^{1/h}$ pairs the $R$-basis $t^{\alpha/h}$ with its
complementary basis $t^{((h-1)\mathbf1-\alpha)/h}$, and has invertible
matrix.  Replacing $dt$ by a volume generator multiplies by units on
the source and target and preserves perfectness.  This establishes
the claim without applying a finite-type smoothness theorem to the
completed power-series ring.

The torsor augmentation has the same dual-generator property.  Here
is a finite-level verification.  The coordinate algebra of the
translation quotient is a tensor product of truncated polynomial
algebras $\kk[z_i]/(z_i^h)$, with its Honda Hopf structure.  A nonzero
translation-invariant functional $\lambda$ gives a pairing
$\lambda(ab)$.  Its radical is a translation-stable ideal: invariance
of $\lambda$ and the antipode identity show that the map to the dual
regular representation is equivariant.  If the radical is proper,
it is contained in the augmentation ideal, since this coordinate
algebra is local.  For a translation-stable ideal $I$ contained there,
apply $\id\otimes\epsilon$ to $\Delta(I)\subset\cO(H)\otimes I$;
the counit identity gives $I=0$.  The radical cannot be the whole
algebra because $\lambda\ne0$.  The pairing is therefore perfect.
After a faithfully flat trivialization this is precisely the pairing
of the torsor functional, so its perfectness descends.

There is a useful precision about the comparison.  The dual module is
rank one over the \emph{source} $R^{1/h}$.  Thus the two functionals
initially differ by
\begin{equation}
 \tau(z)=T_{\eta,h}(u z),\qquad u\in(R^{1/h})^\times,
 \label{eq:trans-source-unit}
\end{equation}
not by an a priori constant multiplying the output.  Both maps are
$G$-equivariant, and the generator property makes $u$ unique, hence
$G$-invariant.  By \Cref{thm:tower-comparison,cor:ind-functions}, its
image after perfected geometric extension and full markings is a
constant on $(\mathcal N_m^r)_{\mathrm{ind}}$.  This extension is
injective on the finite algebraic root ring: power-series expansion
on the analytic annuli is injective in each entry of the whole finite
free basis, remains so after the idempotent summand, and full marking
is faithfully flat.  The coefficient-only
$|\kk|$-Frobenius fixes its finite model, and therefore fixes this
constant.  Consequently $u\in\kk^\times$, with the understood finite
enlargement if the model was defined there.  Since $h$ fixes $\kk$,
\cref{eq:trans-source-unit} is the required scalar comparison.
\end{proof}

\subsection{One transfer, all iterates, and whole-model pole bounds}

When $m>1$ put $h_0=q_0^{m-1}$ as in
\cref{eq:trans-powers}; when $m=1$ choose a power $h_0=p^{\nu_0}>1$
fixing the same finite constants.  In either case set
$\nu_0=\log_p h_0$.

\begin{proposition}[The fixed normalized transfer]
\label[proposition]{prop:trans-cartier}
At the finite marking stage of \Cref{prop:trans-volume}, the operator
\begin{equation}
 S(f)=\frac{\Cartier^{\nu_0}(f\eta)}{\eta}\quad(f\in R)
 \label{eq:trans-S}
\end{equation}
is continuous, $\kk$-linear, and $G$-equivariant.  It satisfies
$S(b^{h_0}f)=bS(f)$.  For $m>1$, the maps on $R$ induced by nonzero
translation augmentations $\mathsf N_{q_0^a}\to\kk$, after the root
identification, are nonzero scalar multiples of $S^a$.  For $m=1$ one
can choose the toral frame so that
\begin{equation}
 S\Fr^{\nu_0}=1.
 \label{eq:trans-height-one-split}
\end{equation}

Extend $S$ to $B'$ by $S(f)=S(ef)$, with image in $eB'$.  There is a
fixed integer $c\geq0$ such that, for every integer $D$,
\begin{equation}
 S(x^{-D}A')\subset x^{-\lceil D/h_0\rceil-c}A'.
 \label{eq:trans-pole-bound}
\end{equation}
Every sufficiently large $D_0$ makes
$\mathcal P=x^{-D_0}A'$ stable under $S$.  Every fixed bounded-pole
lattice, including any compact family of cochains with values in it,
is carried into $\mathcal P$ by all sufficiently large iterates.

The volume, its inverse, the whole trace-pairing matrices and their
inverses have one fixed pole bound $c'$.  In the normalized basis
$b_i^{1/h}$, $h=h_0^a$, their transported matrices have pole losses at
most $c'/h$.  The normalized Cartier maps are adjoint to the inclusions
of root Laurent representatives under the residue pairings with the
transported volumes.
\end{proposition}

\begin{proof}
The formula and \Cref{prop:trans-volume} give equivariance and
$h_0^{-1}$-linearity.  This is $\kk$-linearity because $h_0$ fixes the
finite field.  Continuity follows by grouping finitely many residue
exponents in the $h_0$-basis expansion; multiplication by the fixed
volume and its inverse is continuous on bounded-pole steps.
The intermediate volume factors cancel under composition, so
$S^a(f)=\Cartier^{a\nu_0}(f\eta)/\eta$ for every $a\geq1$.

Suppose first that $m>1$.  Put
$s_q=\prod_iD_i^{q-1}\in\Lambda_q$.  A $\Lambda$-linear functional
$\mathsf N_q\to\kk$ is evaluation at a socle element of $\Lambda_q$.
Its space is the one-dimensional line generated by evaluation at
$s_q$.  Make $U$ pro-$p$.  The action on this line is a continuous
character to $\kk^\times$, a group of order prime to $p$, so it is
trivial.  The maps therefore descend with their natural actions.
Successive coinduction composes the basic socle evaluations into
evaluation at
\[
 \prod_{j=0}^{a-1}s_{q_0}(D_1^{q_0^j},\ldots,D_d^{q_0^j})
       =\prod_iD_i^{q_0^a-1},
\]
which is nonzero in $\Lambda_{q_0^a}$.  Thus these composites do not
vanish.  The coherence of \Cref{prop:trans-roots} identifies their
coefficient maps with the root-twisted composites of
$T_{\eta,h_0}$.  In selfmap notation these are $S^a$.
\Cref{prop:trans-volume} identifies any other nonzero augmentation
with a nonzero constant multiple; the constants introduce no
semilinear twist because all chosen powers fix $\kk$.

Suppose now that $m=1$. We use the order-$p$ torsor
$C=R^{1/p}$ under the fixed group $\mu_p^r$ from
\Cref{prop:trans-volume}. Its character decomposition is strongly
graded by $(\Z/p\Z)^r$, with neutral component $R$ and invertible
character components
\citep[Theorem~12.12 and Remark~12.13]{milne2017}.
The strong grading uses the torsor hypothesis: after a faithfully
flat trivialization, each character component is free of rank one
and multiplication between components is an isomorphism. These
properties descend to $R$.
Zariski locally on $\Spec R$, choose generators $z_i$ of the
fundamental character components. Strong grading makes each $z_i$
a unit, with $b_i=z_i^p\in R^\times$. The monomials $z^\alpha$,
$0\leq\alpha_i<p$, form an $R$-basis of $C$. Taking $p$th powers
shows that the $b^\alpha$ form an $R^p$-basis of $R$, so the $b_i$
are a $p$-basis.

For the standard toral Lie basis,
$\delta_i(z_j)=\delta_{ij}z_j$. \Cref{eq:trans-frame} gives
$\partial_i(b_j)=\delta_{ij}b_j$, so the already constructed volume
is locally
\[
 \eta=d\log b_1\wedge\cdots\wedge d\log b_r.
\]
On overlaps a different choice has $z_i'=c_i z_i$ with
$c_i\in R^\times$, hence $b_i'=c_i^p b_i$ and
$d\log b_i'=d\log b_i$. These local descriptions therefore agree
with the global volume. We retain this standard toral normalization;
the comparison constants multiply the torsor functionals.
The logarithmic Cartier formula gives $\Cartier(\eta)=\eta$.
Consequently, for every chosen $h_0=p^{\nu_0}$ fixing $\kk$,
\[
 S(f^{h_0})
   =\frac{\Cartier^{\nu_0}(f^{h_0}\eta)}{\eta}
   =f\,\frac{\Cartier^{\nu_0}(\eta)}{\eta}
   =f.
\]
This proves \cref{eq:trans-height-one-split} at the existing finite
frame stage. The natural auxiliary actions at higher levels are
retained; no splitting of a higher torsion quotient or further
shrink of $U$ is required. No positive-dimensional local distribution
ring is used in this case.

We next prove the pole bound in the whole model.  Fix an $A$-basis
$b_1,\ldots,b_v$ of $A'$.  The finite \'{e}tale base change makes
$B'^{1/h_0}$ free over $B$ on
\[
 b_i^{1/h_0}x^{\alpha_1/h_0}
     y_1^{\alpha_2/h_0}\cdots y_{r-1}^{\alpha_r/h_0},
 \qquad 0\leq\alpha_j<h_0.
\]
The $B$-linear map $T_{\eta,h_0}$ after multiplication by $e$, with
output included in $B'$, is a fixed finite matrix in these bases.
Its entries belong to $B$ and so have a common finite $x$-pole bound.
For $f\in x^{-D}A'$, its root is a sum of
$x^{-D/h_0}a_i^{1/h_0}b_i^{1/h_0}$, with $a_i\in A$.
Group the exponents in each $a_i^{1/h_0}$ modulo the integral lattice.
The resulting coefficients lie in $A$.  Reducing $-D$ modulo $h_0$
loses at most $\lceil D/h_0\rceil$ integral powers of $x$.
The fixed matrix contributes at most a further $c$, proving
\cref{eq:trans-pole-bound} for all $D$, including negative $D$.
This also proves continuity on compact coefficient lattices.

If $D_{a+1}\leq D_a/h_0+c+1$, then
\[
 D_a\leq h_0^{-a}D+(c+1)\frac{1-h_0^{-a}}{1-h_0^{-1}}.
\]
Choose $D_0>(c+1)/(1-h_0^{-1})$.  The same inequality implies
$\lceil D_0/h_0\rceil+c\leq D_0$ and, for each fixed $D$, eventually
$D_a\leq D_0$.  The estimates apply to the whole lattice, uniformly
over any compact parameter set.  A compact bounded-pole cochain
family uses one such $D$, so the assertion about cochains follows as
well.

For the trace and volume assertion, set
\[
 \theta=d\log x\wedge dy_1\wedge\cdots\wedge dy_{r-1}.
\]
Write $\eta=a_R\theta$ with $a_R\in R^\times$.  Extend its coefficient
to the unit $a=a_R+(1-e)\in B'^\times$ and use $a\theta$ on the whole
algebra.  The trace form of the finite \'{e}tale algebra is perfect;
therefore the matrices
\[
 \bigl(\Tr_{B'/B}(b_i b_j)\bigr)_{ij},\qquad
 \bigl(\Tr_{B'/B}(a b_i b_j)\bigr)_{ij}
\]
and their inverses have entries in $B$.  The multiplication matrices
of $e,a,a^{-1}$ do also.  Taking a maximum bounds all their poles by
one $c'$.  Transport by the root isomorphism replaces every matrix
entry by its $h$th root.  Its pole bound in the original $x$-valuation
is consequently $c'/h$.  The volume at this level is
\[
 \eta_h=a^{1/h}\,d\log x^{1/h}\wedge
                   dy_1^{1/h}\wedge\cdots\wedge dy_{r-1}^{1/h}.
\]
This is transport along an isomorphism, not differential pullback
under an inseparable map.

Finally take a finite Laurent principal-part representative $w$ on
one root grid and a bounded-pole coefficient on the next grid, and
pair them by the coefficient residue of their product times
$\eta_h$, after the finite \'{e}tale trace.  In the logarithmic
$x$-coordinate the residue selects exponent zero; in each ordinary
$y$-coordinate it selects exponent $-1$.  \Cref{eq:trans-cartier-coordinates} therefore gives
\begin{equation}
 \langle w,T_{h,hh_0}(f)\rangle_h
       =\langle \iota w,f\rangle_{hh_0},
 \label{eq:trans-residue-adjoint}
\end{equation}
where $T_{h,hh_0}$ is the $h$-root of $T_{\eta,h_0}$ and $\iota$ is
inclusion of the same Laurent expression on the finer grid.
For instance the ordinary one-variable test is
$h_0i+j=-1$ on both sides; the logarithmic test is $h_0i+j=0$.
The coefficient field is fixed by $h_0$, so taking the selected
coefficient introduces no field twist.  Finite \'{e}tale trace
commutes with this calculation, since its root matrix is the
transported trace matrix.  Multiplication by the volume coefficient
and its inverse cancels as in
\cref{eq:trans-cartier-functional}.  Thus the scalar monomial check
proves \cref{eq:trans-residue-adjoint} in the whole basis and then
after $e$.  Continuity extends the identity whenever a convergent
Laurent expression defines a residue functional.  This is the
adjoint identity used below; no trace for the purely inseparable
root extension is involved.
\end{proof}

\begin{lemma}[One bound for a fixed finite diagram]
\label{lem:trans-diagram-bound}
Fix a finite Yoneda diagram of translation representations, descend
it at a finite marking stage, and fix a cohomological degree.  Its
successive boundary maps can be represented by an operator on
cocycles that loses at most one fixed $x$-pole order.  Under
\Cref{prop:trans-roots}, all of its coinduced diagrams use the roots
of this same operator and its chosen splittings.
\end{lemma}

\begin{proof}
Every term and image bundle is finite locally free on the affine
scheme $\Spec R$.  Each surjection onto such an image has an
underlying $R$-linear section.  Choose these finitely many sections
once.  Embed their finite projective modules in finite free modules
and choose finitely generated whole $A$-lattices; each section is
then a finite matrix over the localization and has a fixed pole
loss.  The group action on a fixed lattice has a uniform pole bound,
by the finite torsor realization and the bounded-step continuity of
\Cref{lem:cts-exact,prop:mark-whole-model}.  Since $g(x)/x$ is an
integral unit, the same bound applies to every shift of that lattice.
The finitely many group differentials and chosen sections computing
the composite boundary thus have a common total pole loss.  This
constructs the required operator on continuous cocycles, with their
bounded-pole convention.

The natural identity \cref{eq:trans-root-functor} carries the entire
diagram, including its arrows, to its root diagram.  Transport each
of the already chosen sections by the same root isomorphism.
Computing the boundary with those sections is exactly the root of
the original computation.  Hence the increasing ranks of the
coinduced representations require no new choices of matrices or
new pole estimate.  This statement concerns realization of a fixed
diagram; comparisons of extension classes after one common
auxiliary shrink are supplied in \Cref{sec:pro-transfer}.
\end{proof}
